% ==========================================================
% 连续命中、命中段与概率
% 难度：★★★ 难题
% ==========================================================

\newcommand{\qProbPermConsecutiveHitProbStem}{%
连续投掷七只飞镖。每次：命中得 $1$ 分；未命中得 $0$ 分；连续命中 $2$ 次额外加 $1$ 分；连续命中 $3$ 次额外加 $2$ 分；以此类推。每次命中概率为 $\frac{2}{3}$，各次投掷相互独立。求总得分恰为 $7$ 分的概率。
}

\newcommand{\qProbPermConsecutiveHitProbAnswer}{%
$\frac{256}{2187}$
}

\newcommand{\qProbPermConsecutiveHitProbSolution}{%
\textbf{第一步：明确计分结构}\\
设一共命中 $H$ 次，形成 $r$ 个最大连续命中段。\\
每个长度为 $k$ 的命中段贡献 $k+(k-1)=2k-1$ 分。\\
因此总得分为 $S=2H-r$。要求 $2H-r=7$。\\
未命中次数为 $7-H$，命中段数最多为 $7-H+1$。\\
检验可得只有两类：$(H,r)=(4,1)$ 或 $(5,3)$。

\textbf{第二步：计算情形一 $(H=4, r=1)$}\\
四次命中必须连续（形成“中中中中”整体）。\\
长度为 $4$ 的连续段可以从第 $1,2,3,4$ 次开始，\\
共有 $4$ 种序列。该情形概率为：\\
\[ P_1 = 4\left(\frac{2}{3}\right)^4\left(\frac{1}{3}\right)^3 = \frac{64}{2187} \]

\textbf{第三步：计算情形二 $(H=5, r=3)$}\\
五次命中要形成三个连续段，\\
两次未命中必须作为分隔。\\
设三个命中段长度为 $x_1,x_2,x_3$，\\
则 $x_1+x_2+x_3=5\ (x_i\geqslant 1)$。\\
其正整数有序解个数为 $C_4^2=6$ 种结构。\\
这对应 $6$ 种合法的命中序列。\\
该情形概率为：
\[ P_2 = 6\left(\frac{2}{3}\right)^5\left(\frac{1}{3}\right)^2 = \frac{192}{2187} \]

\textbf{第四步：得出结论}\\
总概率 $P = P_1+P_2 = \frac{256}{2187}$。
}

\newcommand{\qProbPermConsecutiveHitProbTeachingNotes}{%
\textbf{【避坑与边界说明】}
\begin{itemize}
    \item 连续命中问题不能直接使用普通二项分布，\\
    因为它不关心命中的先后位置和连续情况。
    \item 计分结构“$S=2H-r$”是本题核心。\\
    长度相同的命中段在结构计数中不可重复排列。
    \item 使用隔板法得到的是有序段长：\\
    $(3, 1, 1)$, $(1, 3, 1)$, $(1, 1, 3)$ \\
    被视为不同序列结构。
    \item 连续命中题必须研究：命中段的个数、长度，\\
    以及不命中如何把各段隔开。
\end{itemize}
}
